\section{Recuperation at the label}
\label{sec:15.3}
Every claim a model makes about what's harmful is downstream of labels, and the labels are downstream of a labor arrangement.

An annotator can apply a category well or badly. But what the arrangement gives them no way to do is say the category is wrong — that the taxonomy fails to carve the material, or that the guideline produces the wrong answer here for a reason the guideline can't express. There's no field for that and no time for it, since the quota is set against classification throughput.

So the most informative disagreement is the one the pipeline is least able to receive. Fricker's name for this is hermeneutical injustice: the people best placed to see a wrong lack the shared concepts (here, a field in a form) that would make it intelligible to those with power over them \autocite{fricker2007epistemic}. A misapplied label is a data-quality problem, detectable by inter-annotator agreement and fixable with better instructions.

Inter-annotator agreement is the pipeline's measure of quality, and it therefore works as its selection rule: schemes and annotators that produce agreement are kept, and those that don't are revised or dropped. A selection rule admits only what it can measure. Agreement can measure whether a category was applied consistently; it can't measure whether the category was the right one, because a shared bad scheme produces high agreement as reliably as a good one. That is the metric replacing the thing measured, and exception coded as betrayal is there beside it: the annotator who thinks the scheme is wrong has no field to file that in and no time to file it.

Concept bottleneck models expose named concepts — harm, autonomy, risk — so a person can correct a concept and watch the decision change \autocite{koh2020concept}. What such models also learn is to carry other information \emph{inside} the names, so a coordinate keeps its public meaning and acquires a private one that does the deciding; the obvious mitigations don't remove it \autocite{mahinpei2021promises}. The literature calls it leakage, and sometimes it is. Sometimes the concept is the container that leaks, because the world didn't fit inside it. Purifying the representation protects an explanation from covert information and protects a taxonomy from evidence that its divisions were wrong — same operation, and no measure internal to the taxonomy tells them apart.

\runin{Who is excluded, and from what}

The exclusion is structural, not attitudinal. These people are, for much of the material, the best-placed judges in the pipeline — they've seen more of it than whoever wrote the guideline, and they have the most direct evidence of where the scheme breaks. The arrangement routes their conclusions about individual items upward and their conclusions about the scheme nowhere. And there is an asymmetry: the annotator bears the cost — exposure, injury — and has no stake in the outcome and no claim on it. A worker who can flag an item but can't decline a scheme, and loses the contract by pressing, has a channel and no refusal right.

Pay is a separate obligation: raising it changes the labor arrangement without adding the missing channel.

The checkable question isn't whether a guideline was ever changed because an annotator contested one, since routine revisions happen without changing anyone's standing. The check is what the contest produced — a revision the annotator can point to and argue with again, or a new category in somebody else's taxonomy, a quality score they're now measured against, a reviewer tier they're not in. And ask what happened to the person who raised it.
